\documentclass[../../main.tex]{subfiles}
\begin{document}

{\bf [解答]}

与えられた関数
\begin{align}
  F(x) = \int_{0}^{x} t\sin^3 t dt \label{eq:1}
\end{align}
の極大値を考える．両辺を$x$で微分して
\begin{align}
  F'(x) = x \sin^3 x \label{eq:2}
\end{align}
である．題意より$x>0$だからこの正負は $\sin x$ の正負と等しい．
よって$F'(x)$ は $x = (2n-1)\pi$ ($n \in \mathbb{N}$) の前後で 正から負へと符号がかわるから $x = (2n-1)\pi$ で極大である.

\begin{table}[htb]
  \centering
  \caption{$x = (2n-1)\pi$ 付近における $F(x)$ の増減表 ($n \in \mathbb{N}$)}
  \label{tab:1970_5_variation}
  \renewcommand{\arraystretch}{1.3}
  \begin{tabular}{c|c|c|c|c|c|c|c}
    \hline
    $x$ & $\dots$ & $(2n-2)\pi$ & $\dots$ & $(2n-1)\pi$ & $\dots$ & $2n\pi$ & $\dots$ \\ \hline
    $F'(x)$ & & $0$ & $+$ & $0$ & $-$ & $0$ & \\ \hline
    $F(x)$ & & & $\nearrow$ & 極大 & $\searrow$ & & \\ \hline
  \end{tabular}
\end{table}

一方で$F(x)$を直接計算すると，
\begin{align}
  F(x)
  &= \dfrac{1}{4} \int_0^x t (3\sin t - \sin 3t) \, dt \\
  &= \dfrac{1}{4} \left[ \left( -3t\cos t + 3\sin t \right) - \left( -\dfrac{1}{3}t\cos 3t + \dfrac{\sin 3t}{9} \right) \right]_0^x \\
  &= \dfrac{1}{4} \left[ -3x\cos x + 3\sin x + \dfrac{1}{3}x\cos 3x - \dfrac{1}{9}\sin 3x \right]
\end{align}
だから，$x=(2n-1)\pi$を代入して
\begin{align}
  F((2n-1)\pi) &= \dfrac{1}{4} \left[ +3(2n-1)\pi + \dfrac{1}{3}(2n-1)\pi (-1) \right] \\
  &= \dfrac{2}{3}(2n-1)\pi \quad (n \in \mathbb{N})
\end{align}
が求める極大値である．

\newpage

\end{document}
